% SGA 3, 第 VI_B 篇, 第 1 轮英文重构底本的中译.
% 范围: 命题 8.4 及其完整证明, 包括原星号注及编者注 82--84.
% 权威底本: Polo--Gille Exp6B-13oct24.pdf, 组件第 63--65 页
% (合卷第 404--406 页; 重版页眉第 393--395 页; 原版边栏页码第 377--379 页).
% 第 9 节始于组件第 65 页, 不在本组件内.
% OCR 与 Reinhold 英文 Markdown 仅为对照材料. 本组件无图.

\medskip
\phantomsection\label{prop:sga3-VIB-8.4}%
\noindent\textbf{命题 8.4.---}
\emph{设 \(G\) 是有限呈示的 \(S\)-群,
使得对每个 \(s\in S\), 纤维 \(G_s\) 在 \(\kappa(s)\) 上平滑.
设 \(T\) 是 \(S\) 中使 \(G_s\) 可解\textup{(分别地, 幂零)} 的那些点 \(s\) 所成的集合.}
\begin{enumerate}
\item[\textup{(i)}]\phantomsection\label{concl:sga3-VIB-8.4-i}
  \emph{则 \(T\) 在 \(S\) 中局部可构造.}

\item[\textup{(ii)}]\phantomsection\label{concl:sga3-VIB-8.4-ii}
  \emph{此外, 若假设 \(G\) 在 \(S\) 上平坦且分离
  \textup{(即 \(G\) 在 \(S\) 上平滑、拟紧且分离时)},
  则 \(T\) 在 \(S\) 中是闭的.}
\end{enumerate}

\emph{证明.}
显然可以假设 \(S\) 仿射, 其环为 \(A\).
由 \hyperref[par:sga3-VIB-10.1]{10.1} 及
\hyperref[par:sga3-VIB-10.10]{10.10\textup{b)}},%
\begingroup
\renewcommand{\thefootnote}{\fnsymbol{footnote}}%
\footnote[1]{\label{note:sga3-VIB-8.4-star}
本证明中将使用\hyperref[sec:sga3-VIB-10]{第 10 节} 建立的结果;
与\hyperref[sec:sga3-VIB-9]{第 9 节} 一样,
它们不依赖于当前的\hyperref[sec:sga3-VIB-8]{第 8 节}.}%
\endgroup\ 
存在 \(A\) 的一个 Noether 型子环 \(A'\)
及一个有限型的 \(A'\)-群 \(G'\), 使得
\[
  G'\otimes_{A'}A\simeq G.
\]
由 EGA~IV\(_3\), 11.2.6 和 8.10.5,%
\footnote{\label{note:sga3-VIB-82}编者注: 已将引用 10.8.5 改正为 8.10.5.}
若 \(G\) 在 \(S\) 上平坦且分离,
可以假设 \(G'\) 在 \(S'=\operatorname{Spec}A'\) 上平坦且分离.%
\footnote{\label{note:sga3-VIB-83}编者注: 以下内容作了扩充.}
由于 \(G\) 在 \(S\) 上有限呈示,
\(S\) 中使 \(G_s\) 几何约化,
或等价地在 \(\kappa(s)\) 上平滑的那些点 \(s\) 所成的集合,
是局部可构造的 (EGA~IV\(_3\), 9.7.7).
所以, 由 EGA~IV\(_3\), 9.3.3,
可以假设对每个 \(s'\in S'\), \(G'_{s'}\) 在 \(\kappa(s')\) 上平滑.
另一方面, 若 \(s'\) 表示 \(s\) 在 \(S'\) 中的像, 则
\[
  G'_{s'}\otimes_{\kappa(s')}\kappa(s)\simeq G_s.
\]
由 \hyperref[prop:sga3-VIB-8.3]{8.3},
\(G_s\) 可解 (分别地, 幂零), 当且仅当 \(G'_{s'}\) 也如此.
因此归约到 \(S\) 是 Noether 型仿射概形的情形.

先证明 \(T\) 可构造.
应用 EGA~0\(_{\mathrm{III}}\), 9.2.3 的判别法,
并如前论证, 便归约为证明:
当 \(S\) 为 Noether 型整概形时,
\(T\) 或 \(S-T\) 包含 \(S\) 的一个非空开子集.

于是, 假设 \(S\) 是 Noether 型整概形, 泛点为 \(\eta\).
对每个 \(s\in S\), 置 \(K_s^0=G_s\),
并对 \(p\geq1\) 置
\phantomsection\label{form:sga3-VIB-8.4-Ksp}%
\[
  K_s^p=(K_s^{p-1},K_s^{p-1})
  \qquad
  \bigl(\textup{分别地, }\
    K_s^p=(G_s,K_s^{p-1})\bigr).
\tag*{\(\dagger\)}
\]
\begingroup
\renewcommand{\thefootnote}{\fnsymbol{footnote}}%
\footnotetext[2]{\label{note:sga3-VIB-8.4-G-vs-Gs}
译者注: 此公式中幂零情形的分支, 法文原文印为 \(G\).
由于 \(K_s^{p-1}\) 是一个 \(\kappa(s)\)-群,
且后续公式使用 \(G_s\), 此处恢复了纤维下标.}%
\endgroup
先证明闭子概形的序列 \(K_\eta^p\) 最终稳定.
由 \hyperref[concl:sga3-VIB-7.3-v]{7.3\textup{(v)}},
每个 \(K_\eta^p\) 都正规, 所以序列 \((K_\eta^p)_p\) 递减.
由于 \(G_\eta\) 为 Noether 型, 此序列因而最终稳定.
所以存在一个整数 \(n\), 使得
\phantomsection\label{form:sga3-VIB-8.4-stationary}%
\[
  K_\eta^p=K_\eta^n
  \qquad\text{对每个 }p\geq n.
\]

另一方面, 由 \hyperref[cor:sga3-VIB-10.12.1]{10.12.1}
及 \hyperref[cor:sga3-VIB-10.13]{10.13},
存在 \(S\) 的一个非空开子集 \(S'\)
及一个有限呈示的 \(S'\)-群 \(D\),
使得对每个 \(s\in S'\), 有
\phantomsection\label{form:sga3-VIB-8.4-D-fibers}%
\[
  D_s=K_s^n,
  \qquad
  (D_s,D_s)=D_s
  \quad
  \bigl(\textup{分别地, }(G_s,D_s)=D_s\bigr).
\]
可以假设 \(S'=S\).
于是对每个 \(s\in S\) 及每个 \(p\geq n\), 有 \(D_s=K_s^p\).
所以 \(G_s\) 可解 (分别地, 幂零),
当且仅当 \(D_s\) 同构于单位 \(\kappa(s)\)-群.

由 EGA~IV\(_3\), 9.6.1\textup{(xi)},
\(S\) 中使结构态射
\[
  D_s\longrightarrow\operatorname{Spec}\kappa(s)
\]
为同构的那些点 \(s\) 所成的集合是可构造的.%
\footnote{\label{note:sga3-VIB-84}编者注: 此处使用了 \(S\) 按假设仿射,
因而拟紧且拟分离的事实
\textup{(参见 EGA~0\(_{\mathrm{III}}\), 9.1.12)}.}
因此, 它或者无处稠密, 或者包含 \(S\) 的一个非空开子集.
由此证明 \(T\) 局部可构造.

现在证明, 若 \(G\) 还在 \(S\) 上平坦且分离, 则 \(T\) 是闭的.
由于 \(T\) 局部可构造, 它是闭的当且仅当它在特殊化下稳定
(EGA~IV\(_1\), 1.10.1).

设 \(s\in S\), \(s'\) 是 \(s\) 在 \(S\) 中的一个特殊化.
由于已经归约到 \(S\) 为 Noether 型的情形,
由 EGA~II, 7.1.9 得到一个离散赋值环 \(A\) 及一个态射
\[
  \operatorname{Spec}A\longrightarrow S
\]
使得 \(s\) (分别地, \(s'\)) 是 \(\operatorname{Spec}A\)
的泛点 \(\alpha\) (分别地, 闭点 \(a\)) 的像.
只需证明, 若
\[
  G'=G\times_S\operatorname{Spec}A
\]
且 \(G'_\alpha\) 可解 (分别地, 幂零),
则 \(G'_a\) 也如此.
由于 \(G\) 在 \(S\) 上平坦且分离,
\(A\)-群 \(G'\) 平坦且分离.%
\begingroup
\renewcommand{\thefootnote}{\fnsymbol{footnote}}%
\footnote[3]{\label{note:sga3-VIB-8.4-Gprime-alpha}
译者注: 法文原文印为 \(G'_\alpha\) 在 \(A\) 上平坦且分离.
所指对象应为基变换所得的 \(A\)-群 \(G'\);
其泛纤维 \(G'_\alpha\) 是 \(\kappa(\alpha)\) 上的群.}%
\endgroup\ 
因此归约到 \(S\) 是离散赋值环 \(A\) 的谱的情形.

由于假设 \(G_\alpha\) 可解 (分别地, 幂零),
按以上引入的记号, 存在一个整数 \(q\),
使得 \(K_\alpha^q\) 同构于单位 \(\kappa(\alpha)\)-群.
对每个 \(n\), 以 \(\overline{K_\alpha^n}\) 表示
\(K_\alpha^n\) 在 \(G\) 中按 EGA~IV\(_2\), 2.8.5 意义的概形闭包.
对 \(p\) 作归纳, 证明
\phantomsection\label{form:sga3-VIB-8.4-closure-induction}%
\[
  \bigl(\overline{K_\alpha^p}\bigr)_a
    \supset K_a^p.
\]
首先, 由于 \(G\) 在 \(A\) 上平坦,
由 EGA~IV\(_2\), 2.8.5 有 \(\overline{G_\alpha}=G\), 因而
\[
  \bigl(\overline{K_\alpha^0}\bigr)_a=K_a^0.
\]

设 \(p\geq0\), 并假设
\[
  K_a^p\subset\bigl(\overline{K_\alpha^p}\bigr)_a.
\]
遵循 \hyperref[def:sga3-VIB-7.2.2]{7.2.2},
以 \(\nu_a,\nu_\alpha,\overline{\nu},\overline{\nu}_a\)
表示下列态射. 在可解情形中, 它们为
\phantomsection\label{form:sga3-VIB-8.4-nu-solvable}%
\[
\begin{aligned}
  \nu_a:\;&K_a^p\times_{\kappa(a)}K_a^p
    \longrightarrow G_a,\\
  \nu_\alpha:\;&K_\alpha^p\times_{\kappa(\alpha)}K_\alpha^p
    \longrightarrow G_\alpha,\\
  \overline{\nu}:\;&
    \overline{K_\alpha^p}\times_A\overline{K_\alpha^p}
    \longrightarrow G,\\
  \overline{\nu}_a:\;&
    \bigl(\overline{K_\alpha^p}\bigr)_a
    \times_{\kappa(a)}
    \bigl(\overline{K_\alpha^p}\bigr)_a
    \longrightarrow G_a;
\end{aligned}
\]
分别地, 在幂零情形中, 它们为
\phantomsection\label{form:sga3-VIB-8.4-nu-nilpotent}%
\[
\begin{aligned}
  \nu_a:\;&G_a\times_{\kappa(a)}K_a^p
    \longrightarrow G_a,\\
  \nu_\alpha:\;&G_\alpha\times_{\kappa(\alpha)}K_\alpha^p
    \longrightarrow G_\alpha,\\
  \overline{\nu}:\;&G\times_A\overline{K_\alpha^p}
    \longrightarrow G,\\
  \overline{\nu}_a:\;&
    G_a\times_{\kappa(a)}
    \bigl(\overline{K_\alpha^p}\bigr)_a
    \longrightarrow G_a.
\end{aligned}
\]
由于 \(\nu_\alpha\) 通过 \(K_\alpha^{p+1}\) 分解,
\(\overline{\nu}\) 通过 \(\overline{K_\alpha^{p+1}}\) 分解.
后者显然是 \(G\) 的群子概形,
因而包含 \(\Gamma_G(\overline{\nu})\).
由 \hyperref[concl:sga3-VIB-7.1-iii]{7.1\textup{(iii)}},
\[
  \Gamma_G(\overline{\nu})_a
    =\Gamma_{G_a}(\overline{\nu}_a).
\]
此外, 归纳假设给出
\phantomsection\label{form:sga3-VIB-8.4-induction-inclusions}%
\[
\begin{aligned}
  K_a^p\times_{\kappa(a)}K_a^p
  &\subset
  \bigl(\overline{K_\alpha^p}\bigr)_a
  \times_{\kappa(a)}
  \bigl(\overline{K_\alpha^p}\bigr)_a
  &&\text{在可解情形中},\\
  G_a\times_{\kappa(a)}K_a^p
  &\subset
  G_a\times_{\kappa(a)}
  \bigl(\overline{K_\alpha^p}\bigr)_a
  &&\text{在幂零情形中}.
\end{aligned}
\]
所以,
\phantomsection\label{form:sga3-VIB-8.4-Gamma-chain}%
\[
  K_a^{p+1}
  =\Gamma_{G_a}(\nu_a)
  \subset\Gamma_{G_a}(\overline{\nu}_a)
  =\Gamma_G(\overline{\nu})_a
  \subset\bigl(\overline{K_\alpha^{p+1}}\bigr)_a.
\]

最后, \(K_\alpha^q\) 同构于单位 \(\kappa(\alpha)\)-群.
单位 \(A\)-群在 \(A\) 上平坦,
且同构于 \(G\) 的一个闭子概形,
因为 \(G\) 在 \(A\) 上分离
\textup{(参见 \hyperref[prop:sga3-VIB-5.1]{5.1})}.
由 EGA~IV\(_2\), 2.8.5 可知,
\(\overline{K_\alpha^q}\) 同构于单位 \(A\)-群.
由于已经证明
\[
  K_a^q\subset\bigl(\overline{K_\alpha^q}\bigr)_a,
\]
群 \(K_a^q\) 同构于单位 \(\kappa(a)\)-群.
因此 \(G_a\) 可解\textup{(分别地, 幂零)}.
